\section{Related Work}
\label{sec:related}

\paragraph{Hallucination and its mitigation.}
Hallucination in generation is surveyed broadly by Ji et
al.~\cite{ji2023hallucination} and Huang et al.~\cite{huang2025hallucination},
and profiled specifically for law by Dahl et al.~\cite{dahl2024legalfictions},
whose findings motivate our conservative stance. Retrieval
augmentation~\cite{lewis2020rag, gao2023ragsurvey, karpukhin2020dpr} is the
dominant mitigation; self-consistency detectors such as
SelfCheckGPT~\cite{manakul2023selfcheckgpt} probe the model against itself. Our
contribution is orthogonal and complementary: rather than making the model
\emph{more likely} to be right, we \emph{verify} each citation against an
external source of truth and fail closed otherwise.

\paragraph{Fact verification and legal NLP.}
The FEVER line of work~\cite{thorne2018fever} frames verification as retrieving
evidence and deciding entailment~\cite{williams2018mnli}; we adopt the
retrieve-then-decide shape but specialize the retrieve step to an exact,
state-aware lookup in a versioned legal database, and (for now) the decide step
to existence rather than entailment. In legal NLP, domain encoders and
benchmarks~\cite{chalkidis2020legalbert, guha2023legalbench} target legal
\emph{reasoning}; we instead target legal \emph{citation integrity}, a narrower
but safety-critical property.

\paragraph{Systems and sovereignty.}
On the systems side, the work is a small study in making per-citation checks
cheap over an irregular open-government schema, and in decoupling a trust
pipeline from the serving silicon~\cite{modularmax2024, qualcommcloudai100} via
an OpenAI-compatible indirection. Exposing verification as an MCP
service~\cite{mcp2024} places the guarantee where other agents can reuse it,
turning a single application into shared trust infrastructure.

\paragraph{Attribution and grounded citation.}
A parallel line asks models to \emph{attribute} their statements to sources and
measures whether the attribution holds. We take the strict end of that
spectrum: rather than scoring attribution \emph{post hoc}, we make an
unverifiable attribution \emph{fatal} to the answer. In that sense our closest
relatives are not chatbots but fact-checking and database-integrity systems,
specialized to the narrow, high-stakes surface of statutory citation.
